\chapter{Tool Use and the Agent Loop}
\label{chapter:AgentLoop}
% EVOLVE-BLOCK-START

The loop described in Chapter~1 becomes concrete when the researcher must actually implement it.
Building the loop yourself, rather than configuring a pre-built one, means understanding exactly what the agent can and cannot do, and knowing where to intervene when it fails.
Building a research agent means writing tool definitions, choosing how to structure the control loop, specifying what the agent should return, and anticipating the ways the loop can fail.
This chapter develops each of these in turn.
None of it requires new mathematics, and the program is smaller than it sounds: an agent of the kind built here is a few dozen lines of ordinary Python wrapped around one repeated network call.
Your code sends the conversation so far to the model's API (the network endpoint a program sends text to and reads a reply from), inspects the reply, and, when the reply asks for a tool, runs the matching Python function and sends its result back.
The machinery is small; what is unfamiliar is that one participant in the loop is not deterministic, so the program must be written to tolerate a collaborator that will sometimes do the wrong thing.
Section~\ref{section:ToolDefinitions} covers tool definitions and their schemas; Section~\ref{section:ReAct} develops the ReAct control pattern that interleaves reasoning and action; Section~\ref{section:StructuredOutputs} shows how to constrain what the agent returns; Section~\ref{section:FailureModes} catalogs the ways the loop breaks and how to guard against each; and Section~\ref{section:WorkflowsAndAgents} steps back to when a fixed workflow, rather than an agent, is the better thing to build.

A word on emphasis before the mechanics.
The tools and loop built here put the model in charge of its own execution: it chooses each action at run time.
That is one end of the workflow--agent spectrum of Chapter~\ref{chapter:AgenticParadigm}; the other is a program whose steps are fixed in code, and for many research tasks that fixed end is the better build.
The chapter dwells on the agent end not because it is always right, but because it is the newer and less familiar capability: model-guided execution is what recent models make possible, and it is what a researcher arriving from ordinary scripting has not built before.
We learn it here in full, then choose deliberately where on the spectrum each real system belongs, the question Section~\ref{section:WorkflowsAndAgents} returns to.

\section{Agentic Tool Definitions}
\label{section:ToolDefinitions}

In the context of building agents, a \emph{tool}\index{Tool!definition} is defined by three things: a name, a natural-language description, and a parameter schema.
The model reads all three when deciding which tool to call and how to call it.
The third of these is the part most likely to be unfamiliar: a \emph{parameter schema}\index{Schema} is a machine-readable declaration of the arguments a tool accepts, naming each one, giving its type (text, a number, a list) and marking which are mandatory.
A serviceable mental image is the header of a well-documented function together with its docstring, except that both are handed to the model as data to read rather than compiled into a fixed call site.
Tool definitions are therefore as much a form of documentation as they are a programmatic contract.

This is a genuine departure from how software has classically been built.
A traditional function is called by code that already knows its exact signature; the contract is rigid and the caller's behavior is fixed.
An agentic tool is invoked instead by a model that reads its natural-language description and decides, by interpretation, whether and how to call it.
The description is closer to a specification in plain language than to a type signature: it carries intent the model must construe, and it leaves a latitude that deterministic logic never permitted.
That latitude is the source of the approach's power, since one clear description lets the model handle inputs and situations the author never enumerated.
It also carries a cost.
Interpretation can go wrong where a fixed call site could not, and, more consequentially, an agent that chooses its own tool use tends to \emph{enlarge its own footprint}: the more freedom it has to decide, the more it presumes broad access to files, commands, and the network.
This is why a tool set is a security surface as much as a capability, the reason the least-privilege discipline of Chapter~\ref{chapter:AgenticRig} exists, and it is the tension Section~\ref{section:WorkflowsAndAgents} resolves into a design choice: a workflow bounds the footprint by fixing the steps; an agent widens it by leaving them open.

\subsection{The Tool Schema}

In the Anthropic API, a tool is specified as a JSON object with the following structure:

\begin{lstlisting}[language=json]
{
  "name": "fit_model",
  "description": "Fit a candidate symbolic expression to the training
    data and return the mean squared error on the held-out validation split.
    The expression must be a Python-evaluable string using variables
    x0, x1, ... corresponding to the feature columns.",
  "input_schema": {
    "type": "object",
    "properties": {
      "expression": {
        "type": "string",
        "description": "A Python expression, e.g. 'x0 * x1 + x0**2'"
      }
    },
    "required": ["expression"]
  }
}
\end{lstlisting}

The tool above belongs to the symbolic-regression setting used throughout the course: a candidate formula goes in as text, and a number measuring how badly it fits comes back.
JSON is nothing more than a text format for nested key-value data, the same shape as a Python dictionary, so the object repays a close reading field by field.
The \code{name} is the identifier the model emits when it wants the tool; it comes back verbatim in the model's reply, and the loop uses it to look up which Python function to run.
The \code{input\_schema} declares the arguments in JSON Schema, a small conventional vocabulary for describing data:
\code{"type": "object"} says the arguments arrive as a set of named fields, \code{properties} lists those fields with a type and a description apiece, and \code{required} names the ones the model may not omit.
Note also what the object does \emph{not} contain, which is any code that fits anything:
the definition is the advertisement, and the implementation is a separate Python function the loop invokes when the model asks for the tool by name.

The description field is the most consequential.
It tells the model \emph{when to call the tool} (what preconditions must hold), \emph{what the tool expects} (the semantics of each parameter), and \emph{what the tool returns} (what the model should do with the result).
A vague description produces inconsistent tool-call behavior; a precise description produces reliable output.
The difference is not stylistic.
Shorten the description above to ``fits a model'' and the model has no way to know that the expression must be Python-evaluable, that the variables are named \code{x0}, \code{x1}, and so on rather than by their scientific names, or that the returned number is computed on held-out data and is therefore something to drive down.
Each omission surfaces later as its own class of malformed call that the loop has to catch.

\subsection{Tool Design Principles}

Several principles guide the design of a useful tool set.

\begin{enumerate}
\item \textbf{One tool, one responsibility.}
Each tool should do one thing and do it well.
A tool that both fits a model \emph{and} plots the residuals is harder for the model to invoke correctly than two separate tools.
Smaller, focused tools compose more reliably.

\item \textbf{Tools are safe to call in isolation.}
A tool that has side effects (writes a file, posts to an API, spends tokens) should say so explicitly in its description, and the researcher should consider whether the agent should be able to call it without approval.

\item \textbf{Return only what the model needs.}
A tool that returns a hundred lines of raw data when the model only needs a score fills the context window unnecessarily.
The fitting tool above should return the error, and perhaps a one-line residual summary, not the ten thousand fitted values it computed on the way there.
Post-process tool output to the informative essentials before returning it.

\item \textbf{Return informative errors.}
When a tool fails, its error message becomes the model's next observation.
An error that says \code{"ValueError: shapes (30,) and (40,) not aligned"} enables the model to self-correct; one that says \code{"error"} does not.
\end{enumerate}

These four are not new inventions; they are classical software-engineering discipline, single responsibility, error handling, minimal interfaces, reappearing at the tool boundary.
That is worth pausing on, because it is a specific case of a pattern that runs through the whole course.
Agentic coding contracts out the \emph{typing}, the agent writes the implementation behind each tool, but it does not contract out the \emph{understanding}.
To specify a tool well, and to tell a sound one from a subtly broken one, the researcher still reasons about data structures, execution order, and how errors arise and propagate: the substance of software engineering, not its syntax.
The time spent writing code shrinks; the need to understand computing does not.
Chapter~\ref{chapter:AgenticCoding} made this point for code the researcher owns; it returns here for the tools an agent is given.

\begin{remark}
Tool inputs and results need not be text.
With a vision-capable model, a tool can return an image or a PDF page, and the model reads it directly: a microscopy frame, a gel scan, a plotted residual, or a figure from a retrieved paper becomes an observation like any other.
By this writing in 2026, passing an image alongside text is a routine agent action, and the design principles above carry over unchanged, describe the tool clearly, return only what the model needs, and report informative errors, whatever the modality.
\end{remark}

\subsection{The Tool Call in Practice}

When the model invokes a tool, the loop executes the corresponding function, collects the result, and appends it to the context as a \emph{tool result message}.
The model reads this result alongside the prior context and decides whether to call another tool or produce a final response.

It is worth being concrete about what ``the context'' is in code, because the word suggests something more mysterious than it is.
Your program holds a list of messages, each tagged with a role: what the researcher asked, what the model replied, what each tool returned.
Every turn sends that entire list again, because the model retains nothing between calls, so the growing list \emph{is} the agent's memory and appending a tool result is literally how the agent comes to know what happened.
This is also why the context-window budget of Chapter~\ref{chapter:ContextEngineering} governs how long an agent can run: each step makes the list longer, and the list is resent in full.

\begin{example}
A typical two-step tool use sequence:
\begin{lstlisting}[language=python]
# Turn 1 — model proposes a candidate and calls fit_model
# [model output]:
#   I'll try a quadratic form first.
#   <tool_use name="fit_model">
#     {"expression": "x0**2 + x1"}
#   </tool_use>

# Turn 2 — loop appends the tool result
# [tool_result]: MSE = 0.0023

# Turn 3 — model reads the result and acts
# [model output]:
#   The quadratic fit is very good (MSE = 0.0023).
#   I'll also try x0**2 * x1 to see if an interaction helps.
#   <tool_use name="fit_model">
#     {"expression": "x0**2 * x1"}
#   </tool_use>
\end{lstlisting}
The angle-bracket markers are a reading convention for the page, not characters the model types: in the API the request arrives as an already-parsed field of the reply, carrying a tool name and a dictionary of arguments.
The rhythm is what matters here.
The model never executes anything itself; it asks, and the loop decides whether and how to honor the request, which is precisely where a researcher retains control over an otherwise autonomous process.
\end{example}

\begin{remark}
A single turn need not carry a single tool call.
When a task fans out into independent sub-tasks, such as fitting five candidate expressions or querying three data sources, a frontier model can emit several tool calls in one response, and the loop executes them concurrently, appends all the results as a batch, and re-invokes the model once (Anthropic, 2025).
This \emph{parallel tool calling}\index{Parallel!tool calling} is the idiomatic way to run the fan-out patterns of Chapter~\ref{chapter:Orchestration}: the round trips collapse from one per call to one per batch, and the model reasons over the whole set of results together.
The loop's obligations are unchanged, each result is still validated at its own tool boundary, but its accounting becomes per batch rather than per call.
\end{remark}

\section{The ReAct Pattern}
\label{section:ReAct}

\emph{ReAct}\index{ReAct} (Reasoning + Acting) is a specific formulation of the agent loop that makes the model's intermediate reasoning explicit before each action.
In a ReAct agent, the model alternates between \emph{thought} steps (natural-language reasoning about what to do next) and \emph{action} steps (tool calls or final responses); Figure~\ref{figure:ReAct} traces one such loop from an opening thought to a final answer.

\begin{figure}[htb!]
\begin{center}
%% react.tex: the ReAct thought/action loop (Chapter 6). House conceptual style.
%% Deliberately re-uses the palette of figure:AgentLoop, since ReAct is that loop
%% made concrete: Observation is the Observe step (blue), Thought is Reason (orange),
%% and both a tool call and a final answer are Act (green). The observation is drawn
%% dashed because it is data returned by the world, not a step the model takes.
%% _concept-style.tex: shared TikZ styles for the course's CONCEPTUAL block diagrams.
%% The leading underscore marks this as the one file in figures/ that is not a
%% figure; every other figures/*.tex holds exactly one figure body.
%%
%% \input this at the top of every conceptual figure (before \begin{tikzpicture})
%% so each such diagram speaks one visual language: rounded "block" nodes, stealth
%% arrows, and natural `<hue>!15' fills. Redefining these styles on each \input is
%% idempotent and harmless. The reference for the house parameters is
%% resources/STYLE-figures.md. Figure~\ref{figure:AgentLoop} is the exemplar.
%%
%% Scope: use for conceptual/relational diagrams (boxes-and-arrows). Do NOT use for
%% product mock-ups (the rig figures have their own shared language) or for
%% data/plot figures.
\tikzset{
  % the standard block node: geometry only — apply a `fill=<hue>!15' per node.
  conceptbox/.style={draw, rounded corners=4pt, minimum width=2.4cm,
                     minimum height=0.85cm, align=center, font=\small},
  % a flow arrow between blocks.
  conceptflow/.style={->, >=stealth', thick},
  % an edge/annotation label.
  conceptlabel/.style={font=\scriptsize, align=center},
}%

\begin{tikzpicture}[
    node distance=1.5cm,
    stepbox/.style={conceptbox, minimum width=3.2cm},
    arr/.style={conceptflow},
    ctx/.style={conceptbox, dashed, fill=blue!15,
                minimum width=6.8cm, font=\scriptsize}
]
\node[stepbox, fill=orange!15]   (t1) {Thought};
\node[stepbox, fill=green!15,  right=1.6cm of t1] (a1) {Action (tool call)};
\node[ctx,  below=0.7cm of a1] (ob) {Observation (tool result)};
\node[stepbox, fill=orange!15, below=0.7cm of ob] (t2) {Thought};
\node[stepbox, fill=green!15,  right=1.6cm of t2] (fin) {Final Answer};

\draw[arr] (t1) -- (a1);
\draw[arr] (a1) -- (ob);
\draw[arr] (ob) -- (t2);
\draw[arr] (t2) -- (fin);
\draw[arr] (t2.west) -| (t1.south)
    node[pos=0.35, below, conceptlabel] {more calls};
\end{tikzpicture}%

\caption{The ReAct pattern.
The model alternates between \emph{thought} (explicit reasoning) and \emph{action} (a tool call or final answer).
Each tool result becomes an observation that informs the next thought.
The loop exits when the model produces a final answer rather than another tool call.}
\label{figure:ReAct}
\end{center}
\end{figure}

The explicit thought step serves two purposes.
First, it conditions the subsequent action on intermediate reasoning, reducing errors that arise from the model jumping directly from context to action without an intermediate planning step.
Second, it produces an \emph{audit trail}: by reading the thought steps, the researcher can understand exactly why the agent made each tool call.

The figure is a picture of a program, and the program is short enough to read whole.
Everything the chapter has introduced so far, the tool definitions, the message list, the tool result, meets in about twenty lines.

\begin{example}
The control loop, stripped to its essentials.
\begin{lstlisting}[language=python]
messages = [{"role": "user", "content": task_description}]

for step in range(MAX_STEPS):
    reply = client.messages.create(
        model=MODEL, tools=TOOL_DEFS, messages=messages
    )
    messages.append({"role": "assistant", "content": reply.content})

    calls = [b for b in reply.content if b.type == "tool_use"]
    if not calls:                  # nothing requested: the agent is done
        return reply

    results = []
    for c in calls:                # one turn may request several tools
        try:
            out = TOOLS[c.name](**c.input)
        except Exception as err:   # the error becomes the next observation
            out = f"ERROR {type(err).__name__}: {err}"
        results.append({"type": "tool_result",
                        "tool_use_id": c.id,
                        "content": str(out)})

    messages.append({"role": "user", "content": results})

raise StepBudgetExceeded(MAX_STEPS) # no answer within budget: fail loudly
\end{lstlisting}
\end{example}

Five features of those lines carry the chapter's argument, and each is worth locating in the code.
The \code{messages} list is the agent's whole memory, resent in full on every pass;
\code{TOOL\_DEFS} is the list of JSON objects from Section~\ref{section:ToolDefinitions},
and \code{TOOLS} is the companion dictionary mapping each declared name to the Python function that implements it, which is where the definition and the implementation are finally joined.
The reply is appended \emph{before} any tool runs, so the transcript records what the agent intended even when the tool then fails.
The exit condition is the absence of a tool request: an agent stops when the model answers in prose rather than asking for an action, which means nothing in the loop itself decides that the work is finished.
The \code{try}/\code{except} block is not defensive clutter but the mechanism of Section~\ref{section:FailureModes} in miniature,
since the formatted exception is appended as an observation the model reads on its next pass and can act on, which is why a tool that fails informatively is worth more than one that fails silently.
And the bound on \code{MAX\_STEPS} is the only thing standing between a confused agent and an unbounded run.

One detail is idiomatic rather than conceptual and tends to puzzle first readers: tool results are sent back in a message with the \code{user} role, because from the model's vantage point they are information arriving from outside itself, in the same channel as the original request.
Note too what is absent from the listing.
There is no \code{Thought:} scaffold and no parsing of the model's prose, because tool calls arrive as structured fields rather than as text to be scraped; the ReAct rhythm of Figure~\ref{figure:ReAct} is realized by the message protocol itself, and the next subsection takes up where the reasoning step has gone.

\subsection{Reasoning Models and Implicit ReAct}

The ReAct pattern makes the thought step \emph{explicit}: the model is prompted to write its reasoning into the transcript before each action.
Some frontier \emph{reasoning models}\index{Reasoning model} with extended thinking move this step inward, reasoning internally before emitting a tool call, so the thought becomes latent rather than a visible transcript entry.
This is the same chain-of-thought mechanism introduced in Chapter~\ref{chapter:ContextEngineering}, now performed by the model on its own rather than elicited by a ``think step-by-step'' instruction.
These two placements are complementary rather than exclusive: reasoning can happen \emph{before} the first action, as planning over the initial prompt, or \emph{during} the loop, as a deliberate step that reasons over a freshly returned tool result before the next call.
The latter can be induced even without a reasoning model by exposing a no-op ``think'' tool, a tool with no external effect whose only job is to give the model a scratchpad turn mid-loop, which reportedly helps most in policy-heavy or long sequential-decision settings and little on simple parallel calls (Anthropic, 2025).

A practical consequence is that hand-writing a \code{Thought:}/\code{Action:} scaffold buys little when the model already deliberates before acting; the builder's lever shifts from a prompt template to a \emph{thinking budget}\index{Thinking budget}, a cap on how many reasoning tokens the model may spend before it must answer.
The cost falls on the audit trail.
When reasoning was an explicit transcript entry, the researcher could read \emph{why} each tool was called; when it is latent that record is gone, and even where the internal reasoning is surfaced it is not necessarily a faithful trace of the model's computation, so it cannot be trusted as a complete account of why the agent acted.
The dependable audit trail is therefore the observable one: the sequence of tool calls and their results, which the loop logs regardless of what the model reasoned internally.
This makes the failure-mode discipline of Section~\ref{section:FailureModes} --- validating at each tool boundary and inspecting intermediate outputs --- matter more, not less, as the thought step goes latent.

\section{Structured Outputs}
\label{section:StructuredOutputs}

An agent that produces unstructured text is an endpoint; an agent that produces \emph{structured output}\index{Structured output} is a component that can be composed with other components.
Structured outputs allow the result of one agent to serve as the input of the next without brittle string parsing.
The failure this avoids is worth naming, because it is the ordinary way small research pipelines rot.
If the agent reports ``the best expression I found is \code{x0*x1 + 3}, with an MSE of 0.0023'', the next stage has to recover two values by pattern-matching over English, and that extraction breaks on the day the model writes ``an MSE of roughly 0.0023'' instead.
Structured output shifts the burden from guessing what the text meant to declaring, in advance, which fields must be present and what type each one has.

\subsection{Output Schemas}

The output of the agent's final action should conform to a schema known in advance.
In Python, this is most naturally expressed as a typed dictionary, a dataclass, or a Pydantic model.
Pydantic is a widely used library in which the fields and their types are declared as an ordinary class,
and the declaration doubles as a validator that rejects anything not conforming: a missing field, an extra one, or an \code{mse} that arrived as the text \code{"0.0023"} rather than as a number.

\begin{example}
A structured output for a model discovery result:
\begin{lstlisting}[language=python]
from pydantic import BaseModel

class ModelResult(BaseModel):
    expression: str          # the discovered expression
    mse: float               # mean squared error on the validation split
    complexity: int          # number of nodes in the expression tree
    score: float             # composite score (accuracy + parsimony)
    justification: str       # one-paragraph rationale
\end{lstlisting}
The agent is instructed to produce output conforming to this schema, and the schema is enforced by the calling code.
The \code{complexity} field is the node count of the expression tree, the parsimony measure defined in Chapter~\ref{chapter:Evaluation}, so this component already anticipates the score the evaluation harness will compute.
\end{example}

\subsection{Enforcing Structure}

Two mechanisms enforce structured output.

\begin{enumerate}
\item \textbf{JSON mode / tool-use pattern.}
The model is given a tool called \code{submit\_result} whose input schema is the desired output schema.
The model must call this tool (with correctly typed arguments) to produce its final answer.
The tool-calling layer strongly guides the model toward schema-valid arguments, and some current APIs enforce the schema server-side; where that guarantee is absent, the calling code should still validate the received arguments and re-prompt on any mismatch.

\item \textbf{Response parsing with retries.}
The model is instructed to produce JSON matching a schema, and the calling code parses and validates the output.
If validation fails, the calling code returns the validation error to the model and requests a corrected response.
\end{enumerate}

The first mechanism makes schema-valid output far more likely and is preferable when the API supports it, but in both cases the calling code remains the authoritative validator.
Where only the second mechanism is available, it is a small loop built exactly like the main one: parse the reply, and on failure append the validator's own complaint (\code{"complexity: field required"}) as the next message, which a capable model usually repairs on the first retry.
Budget those retries the way the outer loop is budgeted.
A schema the model cannot satisfy after two or three attempts is usually an ambiguous schema rather than a failing model, and the fix belongs in the field descriptions rather than in more retries.

\section{Failure Modes}
\label{section:FailureModes}

Agents fail.
Understanding the common failure modes enables the researcher to design more robust loops and to recognize failures before they propagate.

\begin{remark}
When a run misbehaves, the diagnostic unit is the \emph{action/observation trace}, the logged sequence of \code{(tool call, arguments, observation)} triples, not the final output.
Walking that trace localizes the fault to one of four places: the wrong tool was chosen, its arguments were malformed, the tool's description was too thin for the model to use it correctly, or a correct observation was misread.
Designing informative errors helps the model recover \emph{during} the loop; reading the trace is how the builder recovers \emph{after} it.
\end{remark}

\subsection{Hallucinated Tool Arguments}

The model may produce tool calls with arguments that look plausible but are incorrect — for example, referencing a column name that does not exist in the dataset, or using a function from a library that is not installed.
These calls fail at execution and return error messages; the model must recover by correcting the argument or choosing a different action.

Good tool descriptions reduce the frequency of hallucinated arguments.
Tool schemas with \code{enum} constraints for categorical parameters are particularly effective: the model is restricted to valid choices by the schema itself.
An \code{enum} simply lists the permitted values of a field, so a \code{split} parameter declared as one of \code{"train"}, \code{"validation"}, \code{"test"} cannot be called with a plausible-sounding \code{"holdout"}.
The general move is to push whatever can be checked into the schema, where it is enforced before the call, rather than into the description, where it is merely suggested.

\subsection{Loops Without Progress}

The agent loop can enter a cycle: the model calls a tool, receives an error, attempts to fix it, receives the same error, and repeats.
A loop that does not progress must be detected and terminated.
The most common causes are:
\begin{itemize}
\item a mismatch between the tool's behavior and the model's mental model of it,
\item an error message that the model consistently misinterprets, or
\item a constraint in the task that the model cannot satisfy with the available tools.
\end{itemize}
The third cause is the one to watch for, because it looks like stubbornness and is actually an incomplete tool set.
The model calls the fitting tool with a column name it guessed from the task description, the tool answers with a key error, the model tries a near-synonym, and the cycle continues, because nothing in the tool set ever lets it ask what the columns are actually called.
The fix is a tool that lists the schema of the loaded data, not a sterner instruction.
The loop should have a maximum step count; if it is exceeded without producing a final answer, the loop should exit and report the failure.

\subsection{Context Overrun}

If the context window fills before the agent produces a final answer, the model cannot proceed.
The loop must either truncate the history (losing information) or terminate.
Long tool results, such as stack traces, raw data dumps, or verbose API responses, are the most common cause.
The arithmetic is unforgiving: a single failed command can return a stack trace of several thousand tokens, and a handful of those will crowd out the task description the agent is still working from.
Truncating tool results to their informative essentials is the primary mitigation.
For a stack trace, keeping the first and last few lines usually preserves everything diagnostic, since that is where the exception type and the failing call site live.

\subsection{Cascading Errors}

A mistake in an early step propagates silently through subsequent steps.
The model fits a model to the wrong data, reports a score on that data, and selects the model as the best, but the score is meaningless because the data was wrong from the start.
The underlying pattern is familiar from ordinary analysis: a units mismatch in a loaded column, a date filter that silently drops the first year of a panel, a merge that quietly duplicates every matched row.
What is new is only that no one is looking at the intermediate output unless the loop was written to look, and that a fluent final report can make a corrupted run read like a successful one.
Defense against cascading errors requires:
\begin{itemize}
\item validation at each tool boundary (verify that the data loaded correctly before fitting),
\item explicit checks in the loop for known failure conditions (is the validation set non-empty?), and
\item inspection of intermediate outputs, not just final results.
\end{itemize}

\section{Workflows, Agents, and What to Build}
\label{section:WorkflowsAndAgents}

The loop this chapter builds sits at one end of a spectrum.
Chapter~\ref{chapter:AgenticParadigm} drew the distinction between an \emph{agent}, which lets the model direct its own process step by step, and a \emph{workflow}, which routes execution through paths fixed in code (Figure~\ref{figure:WorkflowAgentSpectrum}).
Everything so far, tool definitions, the ReAct loop, structured outputs, has developed the agent end, where the model chooses each action at run time.
The other end is a plain program: the steps are settled in advance, and no model decides anything while it runs.
The point worth drawing out here is that agentic development builds \emph{both}.
The skill is the same; only the build target moves along the spectrum, from model-guided execution to steps fixed in code.

There is real value in aiming, deliberately, at the fixed-code end.
A researcher can use an agent to \emph{write} a deterministic program for a task that is well understood and repetitive, then run that program with no model in the loop at all, or with only a light agentic flow that invokes it.
The agent is the author (Chapter~\ref{chapter:AgenticCoding}); the artifact it leaves behind is an ordinary script.
A repetitive step that can be coded away, a nightly reformat of instrument output, or a fixed cleaning-and-normalization pass over each new batch of assays, is often best turned into exactly such a program: agentically built once, then run unattended for as long as the task holds.

The fixed-code end buys a specific set of properties.
A workflow is \emph{deterministic} (the same input yields the same output), cheap (it spends no tokens per run), fast, and auditable, since the code is its own specification.
It is also \emph{robust to model evolution}: a fixed program does not change when the underlying model is updated, retired, or repriced, whereas an agent's behavior can shift with every model release.
For a pipeline a lab expects to rerun for years, that stability is not a minor convenience; it is part of what keeps results comparable across time.

The agent end buys the opposite properties.
It handles inputs the author did not foresee, adapts its path when the task is open-ended, and needs far less specification up front.
Those gains are paid for in the coin of Section~\ref{section:FailureModes}: nondeterminism, per-run cost and latency, behavior that drifts as models change, and a thinner audit trail.
Neither end is better in the abstract; which to build is a property of the task, not a matter of taste.

Most real systems live between the extremes.
A light agentic flow calls deterministic sub-programs, each fixed tool a small workflow the model merely orchestrates; the tool-design principles of Section~\ref{section:ToolDefinitions} read, in this light, as advice for building the workflow components an agent composes.
There is also a natural lifecycle along the spectrum: explore an unfamiliar task with an agent, and once its procedure stabilizes, compile the settled steps down into code, letting the agent discover the recipe and then fixing it in a program.
Chapter~\ref{chapter:Orchestration} takes up composing such components into reliable systems.

\begin{principle}
\label{principle:FixWhatIsStable}
Fix in code what is stable; reserve agency for what must adapt.
Reach for the workflow end when a task is well understood, repetitive, and expected to recur; reserve full model-guided execution for the genuinely open-ended, and for the parts of a system that must handle what you did not anticipate.
The strongest designs are often hybrids, an agent composing deterministic components, with proven agentic procedures graduating into workflows over time.
\end{principle}

\begin{remark}
A second axis now cuts across the workflow--agent spectrum: how much computation to spend at inference.
\emph{Inference-time compute scaling}\index{Inference-time compute} is the observation that the same model, given a larger \emph{thinking budget} (Section~\ref{section:ReAct}), can solve harder problems, so capability can be bought at run time rather than only by training a larger model.
This decouples capability from model size and adds a knob to the build decision: in an evolutionary search, for instance, a fast, low-budget call proposes routine candidate variants each round, while a slow, high-budget reasoning pass is reserved for the rounds where no candidate has improved in many generations and the search needs a genuinely different idea.
It also cuts both ways for the choice above: it sharpens the case for an agent on open-ended work, where spending more thinking pays, and it raises the running cost a workflow sidesteps, since a fixed program that calls no model spends no inference budget at all.
Treating the thinking budget as an explicit parameter, rather than a fixed property of the model, is part of building the cost-aware systems of Chapter~\ref{chapter:EvolutionaryFrameworks}.
\end{remark}

\section*{Further Reading}

\begin{small}
\begin{enumerate}
\item Yao, S., et al., ``ReAct: Synergizing Reasoning and Acting in Language Models,'' \emph{ICLR}, 2023 (arXiv:2210.03629) — the original formulation of the thought–action–observation pattern.
\item Anthropic, \emph{Tool Use Guide} — the Anthropic API reference for tool definitions, tool-use messages, and enforcing structured outputs.
\item Shinn, N., et al., ``Reflexion: Language Agents with Verbal Reinforcement Learning,'' \emph{NeurIPS}, 2023 (arXiv:2303.11366) — on agents that detect and recover from their own failures.
\item Anthropic, ``Writing Effective Tools for Agents --- with Agents,'' \emph{Anthropic Engineering}, 2025 — a tool definition and its returned result are both prompt surfaces; namespace related tools and return high-signal, human-readable output.
\item Anthropic, ``The `Think' Tool: Enabling Claude to Stop and Think in Complex Tool Use Situations,'' \emph{Anthropic Engineering}, 2025 — a no-op scratchpad tool that adds a mid-loop reasoning turn, with a per-domain evaluation of when it helps.
\item Anthropic, ``Introducing Advanced Tool Use on the Claude Developer Platform,'' \emph{Anthropic Engineering}, 2025 — tool definitions as a context cost, on-demand tool discovery, and programmatic tool calling.
\end{enumerate}
\end{small}
% EVOLVE-BLOCK-END
